\typeout{NT FILE chapter7.tex}%

\chapter{OFLAT Design and Features}
\label{cha:oflat-design}

This chapter deepens the description of the existing \textbf{graphical} OFLAT application (web-based) with a focus 
on how finite-state transducer (FST)
support was integrated. The goal is to show how users interact with FSTs without disturbing established workflows
for other automata models. Implementation details appear in \autoref{cha:oflat-implementation}.

\section{Design Goals}
\label{sec:oflatdesign_goals}
The interface is shaped by four principles:
\begin{itemize}
  \item \textbf{Discoverable actions:} core tasks (add states, create transitions, run, step, inspect) stay on-canvas
	  through direct manipulation and context menus.
  \item \textbf{Pedagogical feedback:} color, overlays, and tables make nondeterminism and acceptance paths explicit;
	  configuration counts and traces match textbook views for FA and extend cleanly to FST (with output words).
  \item \textbf{Consistency:} the same gestures and menus operate for FA, PDA, TM, and now FST;
  \item \textbf{Low-friction:} edits apply in place with immediate visual updates; auxiliary panes (tables, traces)
	  live adjacent to the canvas so users can alternate between graph and data quickly.
\end{itemize}

\section{Workspace and Graph Interaction}
\label{sec:oflatdesign_workspace}
The canvas (Cytoscape.js) remains the primary surface for model construction and inspection. For FSTs, edges carry
      \textit{input:output} label:
\begin{figure}[H]
  \centering
  \includegraphics[width=0.95\textwidth]{fst_labels.png}
  \caption{OFLAT canvas with input:output edge labels.}
  \label{fig:oflat-canvas}
\end{figure}
\begin{itemize}
  \item direct add/move/rename of states; toggling initial/accepting via node menus;
  \item creation and removal of transitions with inline editors that accept both input and output symbols (epsilon allowed);
  \item zoom/pan with popper overlays that stay anchored to nodes even as the viewport moves.
\end{itemize}
Visual coding distinguishes initial and accepting states; a neutral base style is restored before each run so execution
highlights remain readable.

\section{Information Box, Poppers, and Context Menus}
\label{sec:oflatdesign_infobox}
Node context menus expose per-configuration entries during step-by-step execution. Poppers display the count of active
configurations on each node, giving quick feedback on nondeterministic branching. The information box surfaces the current
configuration details: state, remaining input, and accumulated output (for FSTs). All overlays are lightweight and follow
node movement to avoid obscuring the graph.
\begin{figure}[H]
  \centering
  \includegraphics[width=0.95\textwidth]{fst_poppers.png}
  \caption{Information box and popper displaying state/input/output and configuration counts.}
  \label{fig:oflat-infobox}
\end{figure}

\section{Step-by-Step Execution Experience}
\label{sec:oflatdesign_step}
Execution highlights every active state at the current step. In nondeterministic scenarios, the path chosen by the engine
is painted in a darker accent while alternative branches keep a lighter shade, making the selected witness path evident
without hiding other possibilities. Users advance or rewind one step at a time; the input word view marks consumed vs.
remaining symbols with a cursor separator. For FSTs, the accumulated output is shown alongside the remaining input.
\begin{figure}[H]
  \centering
  \includegraphics[width=0.95\textwidth]{fst_step.png}
  \caption{Step-by-step execution with highlighted path, input cursor, and accumulated output.}
  \label{fig:oflat-step}
\end{figure}

\section{Trace Tables and Minimization Views}
\label{sec:oflatdesign_tables}
Tabular views mirror formal definitions: transition tables per input symbol (and input:output for FSTs), execution trace
tables with state, remaining input, and produced output, and minimization colorings that render equivalence classes on the
graph. Tables rebuild automatically when the model changes, keeping visualization and data in sync.
\begin{figure}[H]
  \centering
  \includegraphics[width=0.95\textwidth]{tracefst.png}
  \caption{Trace table for an FST run showing state, remaining input, and produced output.}
  \label{fig:oflat-trace}
\end{figure}

\section{Supported Models and Scope}
\label{sec:oflatdesign_models}
The interaction model is uniform across finite automata, pushdown automata, Turing machines, and finite-state transducers.
For FSTs, edge labels include input/output, configuration displays carry both remaining input and output, and popper counts
reflect configuration multiplicity.

\section{FST-Specific Workflows}
\label{sec:oflatdesign_fstworkflows}
The UI surfaces FST capabilities without diverging from existing gestures:
\begin{itemize}
  \item \textbf{Create and edit transitions:} inline editors accept \texttt{input:output}, with epsilon allowed on either side; 
  validation reuses the library alphabets to flag typos early.
  \item \textbf{Run vs. check outputs:} For exercises, \texttt{acceptExpectedOutput} lets inside sets specify \texttt{input->output} pairs, ensuring 
  exact output matching.
  \item \textbf{Visualization hooks:} the configuration trail from \texttt{acceptFull} drives node highlighting and popper counts 
  (via \texttt{inputEdges}); accumulated output is shown in the information box and trace tables at each step.
  \item \textbf{Transformations and minimization:} determinization/minimization (when applicable) and 
  composition/union/intersection/inverse/concatenation are callable from the same menu used for FA/PDA/TM, 
  returning a new FST rendered on the canvas.
  \item \textbf{Generation:} bounded word generation lists input/output pairs, mirroring the FA/PDA generation UI but including produced 
  outputs.
\end{itemize}

\section{Navigation to the Implementation View}
\label{sec:oflatdesign_nav}
Implementation details, module structure, event wiring, rendering logic, and the bridge from the library to the UI, are
covered in \autoref{cha:oflat-implementation}.